\documentclass[11pt]{article}

% ---- theme variables (passed via pandoc -V) ----
\newcommand{\ThemeName}{Behavioral \& Leadership}
\newcommand{\ThemeKicker}{Soft-Skills Track}
\newcommand{\ThemeNumber}{11}

\usepackage{interviewstyle}
\setthemecolor{9333ea}{c084fc}

% pandoc-required plumbing
\providecommand{\tightlist}{\setlength{\itemsep}{1pt}\setlength{\parskip}{0pt}}
\setlength{\emergencystretch}{3em}
\providecommand{\pandocbounded}[1]{#1}

% syntax-highlighting macros emitted by pandoc (skylighting)
\usepackage{color}
\usepackage{fancyvrb}
\newcommand{\VerbBar}{|}
\newcommand{\VERB}{\Verb[commandchars=\\\{\}]}
\DefineVerbatimEnvironment{Highlighting}{Verbatim}{commandchars=\\\{\}}
% Add ',fontsize=\small' for more characters per line
\usepackage{framed}
\definecolor{shadecolor}{RGB}{35,38,41}
\newenvironment{Shaded}{\begin{snugshade}}{\end{snugshade}}
\newcommand{\AlertTok}[1]{\textcolor[rgb]{0.58,0.85,0.30}{\textbf{\colorbox[rgb]{0.30,0.12,0.14}{#1}}}}
\newcommand{\AnnotationTok}[1]{\textcolor[rgb]{0.25,0.50,0.35}{#1}}
\newcommand{\AttributeTok}[1]{\textcolor[rgb]{0.16,0.50,0.73}{#1}}
\newcommand{\BaseNTok}[1]{\textcolor[rgb]{0.96,0.45,0.00}{#1}}
\newcommand{\BuiltInTok}[1]{\textcolor[rgb]{0.50,0.55,0.55}{#1}}
\newcommand{\CharTok}[1]{\textcolor[rgb]{0.24,0.68,0.91}{#1}}
\newcommand{\CommentTok}[1]{\textcolor[rgb]{0.48,0.49,0.49}{#1}}
\newcommand{\CommentVarTok}[1]{\textcolor[rgb]{0.50,0.55,0.55}{#1}}
\newcommand{\ConstantTok}[1]{\textcolor[rgb]{0.15,0.68,0.68}{\textbf{#1}}}
\newcommand{\ControlFlowTok}[1]{\textcolor[rgb]{0.99,0.74,0.29}{\textbf{#1}}}
\newcommand{\DataTypeTok}[1]{\textcolor[rgb]{0.16,0.50,0.73}{#1}}
\newcommand{\DecValTok}[1]{\textcolor[rgb]{0.96,0.45,0.00}{#1}}
\newcommand{\DocumentationTok}[1]{\textcolor[rgb]{0.64,0.20,0.25}{#1}}
\newcommand{\ErrorTok}[1]{\textcolor[rgb]{0.85,0.27,0.33}{\underline{#1}}}
\newcommand{\ExtensionTok}[1]{\textcolor[rgb]{0.00,0.60,1.00}{\textbf{#1}}}
\newcommand{\FloatTok}[1]{\textcolor[rgb]{0.96,0.45,0.00}{#1}}
\newcommand{\FunctionTok}[1]{\textcolor[rgb]{0.56,0.27,0.68}{#1}}
\newcommand{\ImportTok}[1]{\textcolor[rgb]{0.15,0.68,0.38}{#1}}
\newcommand{\InformationTok}[1]{\textcolor[rgb]{0.77,0.36,0.00}{#1}}
\newcommand{\KeywordTok}[1]{\textcolor[rgb]{0.81,0.81,0.76}{\textbf{#1}}}
\newcommand{\NormalTok}[1]{\textcolor[rgb]{0.81,0.81,0.76}{#1}}
\newcommand{\OperatorTok}[1]{\textcolor[rgb]{0.81,0.81,0.76}{#1}}
\newcommand{\OtherTok}[1]{\textcolor[rgb]{0.15,0.68,0.38}{#1}}
\newcommand{\PreprocessorTok}[1]{\textcolor[rgb]{0.15,0.68,0.38}{#1}}
\newcommand{\RegionMarkerTok}[1]{\textcolor[rgb]{0.16,0.50,0.73}{\colorbox[rgb]{0.08,0.19,0.26}{#1}}}
\newcommand{\SpecialCharTok}[1]{\textcolor[rgb]{0.24,0.68,0.91}{#1}}
\newcommand{\SpecialStringTok}[1]{\textcolor[rgb]{0.85,0.27,0.33}{#1}}
\newcommand{\StringTok}[1]{\textcolor[rgb]{0.96,0.31,0.31}{#1}}
\newcommand{\VariableTok}[1]{\textcolor[rgb]{0.15,0.68,0.68}{#1}}
\newcommand{\VerbatimStringTok}[1]{\textcolor[rgb]{0.85,0.27,0.33}{#1}}
\newcommand{\WarningTok}[1]{\textcolor[rgb]{0.85,0.27,0.33}{#1}}

\begin{document}

\makecoverpage

\thispagestyle{empty}
{\hypersetup{hidelinks}\tableofcontents}
\clearpage

\begin{quote}
\textbf{Audience:} Mid-to-senior engineers preparing for FAANG behavioral rounds.
\textbf{Goal:} Deep mastery --- templates, story banks, worked examples, mnemonics --- ready for prep and night-before revision.
\end{quote}

\section{Overview --- What it is}\label{overview--what-it-is}

Behavioral interviews assess \emph{how} you work, not just \emph{what} you can build. Interviewers ask you to recall real past situations and narrate them. The premise: \textbf{past behavior is the best predictor of future behavior.}

Every major tech company runs behavioral rounds under slightly different branding:

\begin{longtable}[]{@{}
  >{\raggedright\arraybackslash}p{(\columnwidth - 4\tabcolsep) * \real{0.0914}}
  >{\raggedright\arraybackslash}p{(\columnwidth - 4\tabcolsep) * \real{0.4154}}
  >{\raggedright\arraybackslash}p{(\columnwidth - 4\tabcolsep) * \real{0.4931}}@{}}
\toprule\noalign{}
\rowcolor{acc!16}
\begin{minipage}[b]{\linewidth}\raggedright
Company
\end{minipage} & \begin{minipage}[b]{\linewidth}\raggedright
Label
\end{minipage} & \begin{minipage}[b]{\linewidth}\raggedright
Format
\end{minipage} \\
\midrule\noalign{}
\endhead
\bottomrule\noalign{}
\endlastfoot
Amazon & Leadership Principles (LPs) & Every loop interview; Bar Raiser owns the LP rubric \\
\rowcolor{zebra}
Google & Googleyness \& Leadership & 1-2 dedicated rounds; also woven into design/coding debrief \\
Meta & "Move Fast, Be Direct, Build Social Value" & Behavioral + cross-functional fit \\
\rowcolor{zebra}
Microsoft & "Growth Mindset" / Values & Usually 1-2 rounds; culture-fit emphasis \\
Apple & "Apple Values" & Less structured but heavily culture-oriented \\
\end{longtable}

These rounds are \textbf{as eliminatory as coding.} A strong algorithm answer will not save you from a failing behavioral score.

\section{Why It Exists}\label{why-it-exists}

\subsection{What companies are actually buying}\label{what-companies-are-actually-buying}

When a company hires a senior engineer, they are not buying code --- they are buying \textbf{judgment, influence, and force multiplication.} Behavioral rounds exist to answer:

\begin{enumerate}
\def\labelenumi{\arabic{enumi}.}
\tightlist
\item
  \textbf{Will this person make good decisions under pressure and ambiguity?}
\item
  \textbf{Will they own outcomes, or find excuses?}
\item
  \textbf{Can they collaborate without a manager enforcing every handoff?}
\item
  \textbf{Will they raise the bar or lower it?}
\item
  \textbf{Have they done this before at the level we need?}
\end{enumerate}

\subsection{The bar-raiser model (Amazon) and calibrated loops (Google)}\label{the-bar-raiser-model-amazon-and-calibrated-loops-google}

At Amazon, one "Bar Raiser" --- a trained, senior, cross-org interviewer --- has veto power over any hire. Their job is exclusively to assess LPs and ensure the candidate is "above the median of everyone in this role." At Google, a hiring committee reads all scorecards and calibrates against the job ladder. \textbf{Your behavioral answers are being pattern-matched against a written rubric, not vibes.}

\subsection{What gets screened out}\label{what-gets-screened-out}

\begin{itemize}
\tightlist
\item
  Engineers who blame teammates for failures
\item
  People who "we\textquotesingle d"-everything (can\textquotesingle t isolate their own contribution)
\item
  Candidates who have no examples of taking initiative beyond assigned scope
\item
  People who struggle with conflict --- either avoid it or escalate it badly
\item
  Engineers who can\textquotesingle t measure their own impact
\end{itemize}

\section{Why FAANG Cares}\label{why-faang-cares}

\subsection{Amazon --- Leadership Principles \& Bar Raiser}\label{amazon--leadership-principles--bar-raiser}

Amazon codified 16 (now expanded from 14) LPs that define every hiring, promotion, and performance decision. The Bar Raiser is specifically trained to probe LPs. \textbf{Each interview loop is pre-assigned specific LPs} --- you may face 2-3 LPs per interviewer. Failing on any LP can be a veto. Amazon is arguably the most rigorous behavioral interviewer in the industry.

\textbf{Key Amazon signal:} Specificity. Vague answers score 0. They want \emph{exact} situations, \emph{exact} data, \emph{exact} decisions you made personally.

\subsection{Google --- Googleyness \& Leadership}\label{google--googleyness--leadership}

Google\textquotesingle s rubric has two components:

\begin{itemize}
\tightlist
\item
  \textbf{Googleyness:} Comfort with ambiguity, bias to action, humble+confident, collaborative, principled (does the right thing even when hard).
\item
  \textbf{Leadership:} Emergent leadership (took ownership without being asked), influence, driving outcomes.
\end{itemize}

Google is looking for the \textbf{"brilliant friend"} --- someone smart enough to be right but humble enough to learn. Arrogance or "I\textquotesingle m always right" patterns are immediate red flags. So is passivity.

\textbf{Key Google signal:} Ownership + intellectual humility. Did you push something forward while bringing others along?

\subsection{Meta --- "Move Fast, Be Direct"}\label{meta--move-fast-be-direct}

Meta\textquotesingle s behavioral lens centers on:

\begin{itemize}
\tightlist
\item
  \textbf{Moving fast:} Bias to action, ships things, doesn\textquotesingle t over-engineer.
\item
  \textbf{Be Direct:} Constructive candor, disagree openly rather than passive-aggressively.
\item
  \textbf{Build social value:} Impact at scale, cares about the mission.
\end{itemize}

Meta also weights \textbf{cross-functional collaboration} heavily --- did you work well with PM, design, data science?

\subsection{Microsoft --- Growth Mindset}\label{microsoft--growth-mindset}

Microsoft post-Nadella is heavily Dweck-influenced. They want to see:

\begin{itemize}
\tightlist
\item
  Curiosity and learning from failure (not just "I failed and then succeeded" but "here\textquotesingle s what I \emph{learned}").
\item
  Collaboration over competition.
\item
  Empowering others.
\end{itemize}

\subsection{How each company weights behavioral vs. technical}\label{how-each-company-weights-behavioral-vs-technical}

\begin{longtable}[]{@{}
  >{\raggedright\arraybackslash}p{(\columnwidth - 4\tabcolsep) * \real{0.1189}}
  >{\raggedright\arraybackslash}p{(\columnwidth - 4\tabcolsep) * \real{0.2583}}
  >{\raggedright\arraybackslash}p{(\columnwidth - 4\tabcolsep) * \real{0.6229}}@{}}
\toprule\noalign{}
\rowcolor{acc!16}
\begin{minipage}[b]{\linewidth}\raggedright
Company
\end{minipage} & \begin{minipage}[b]{\linewidth}\raggedright
Behavioral Weight
\end{minipage} & \begin{minipage}[b]{\linewidth}\raggedright
Can behavioral override strong technical?
\end{minipage} \\
\midrule\noalign{}
\endhead
\bottomrule\noalign{}
\endlastfoot
Amazon & Very high (\textasciitilde40\%) & Yes --- Bar Raiser LP veto is absolute \\
\rowcolor{zebra}
Google & High (\textasciitilde30\%) & Yes --- hiring committee can downlevel or reject \\
Meta & Medium-high & Yes --- "culture fit" can block \\
\rowcolor{zebra}
Microsoft & Medium & Usually softens a borderline tech score \\
\end{longtable}

\textbf{Takeaway: Behavioral is not a box-check. It is a co-equal signal.}

\section{The STAR Method}\label{the-star-method}

\visualfigure{../../figures/11-behavioral/01-star-method.pdf}{Structure every behavioral story as Situation $\rightarrow$ Task $\rightarrow$ Action $\rightarrow$ Result. Lead with context, but spend most words on your specific actions and a measurable result.}

\subsection{Structure}\label{structure}

\setcodelang{}

\begin{verbatim}
S — Situation   (10-15% of answer time)
T — Task        (5-10%)
A — Action      (60-70%)  ← this is where you score points
R — Result      (15-20%)  ← must be concrete and quantified
\end{verbatim}

\subsection{The STAR principle}\label{the-star-principle}

\begin{itemize}
\tightlist
\item
  \textbf{S:} Set the scene briefly. One or two sentences. Context only.
\item
  \textbf{T:} What was YOUR specific responsibility? (Not the team\textquotesingle s.)
\item
  \textbf{A:} What did YOU specifically do? Be granular. Walk through decisions, tradeoffs, why you chose this over alternatives.
\item
  \textbf{R:} What happened? Numbers. Business outcome. What would have happened without you?
\end{itemize}

\subsection{The Weak vs. Strong worked example}\label{the-weak-vs-strong-worked-example}

\textbf{Question:} "Tell me about a time you had to influence a decision without having authority."

\textbf{WEAK ANSWER:}

\begin{quote}
"We were working on a project and the team wanted to use technology X but I thought technology Y was better. I presented my case and eventually they agreed and we switched. It worked out well and the project was a success."
\end{quote}

\textbf{Why it fails:}

\begin{itemize}
\tightlist
\item
  No specific situation or timeline
\item
  No stakes --- why does it matter which technology?
\item
  "I presented my case" --- what did that mean? One email? A deck?
\item
  "It worked out well" --- zero quantification
\item
  "We switched" --- passive; who made the decision? What convinced them?
\item
  Interviewer learns nothing about your judgment or influence skills
\end{itemize}

\textbf{STRONG ANSWER:}

\begin{quote}
"It was Q3 last year. We were rebuilding our real-time notification service that handled about 2 million events per day. The team lead had already committed to using Kafka for the message queue --- it was on the roadmap and one engineer had prior experience with it. I was not the tech lead, but I had just come off a project where we\textquotesingle d used Kafka at similar scale and hit serious operational complexity around partition rebalancing under traffic spikes.

My task was to make sure we made the right technical decision, even though the default was already set and re-opening it would create friction.

I started by spending two days writing a structured comparison doc --- not a "Kafka is bad" doc, but a "here are the three properties that matter for our use case" doc. I focused on: operational burden on a team with no existing Kafka expertise, latency p99 under our specific write patterns, and vendor lock-in. For each property I used numbers from our past incident data and publicly available benchmarks. I also reached out to the engineer who championed Kafka and had a 1:1 first --- I wanted to understand what they were optimizing for, not argue with them. It turned out they cared most about throughput, not latency. That reframing was key.

I presented the doc in our next architecture review. I explicitly said: "I\textquotesingle m not saying Kafka is wrong. I\textquotesingle m saying here are the specific conditions under which it\textquotesingle s the wrong choice for us." The team lead asked me to run a small proof-of-concept with the alternative (SQS + SNS fanout) over one sprint.

The PoC showed a 40\% reduction in setup complexity and our p99 latency stayed under 200ms at 3x our expected load. We shipped with SQS/SNS. Six months later, a comparable team who did use Kafka hit a 4-hour outage from partition imbalance --- the exact failure mode I had flagged. Our system had zero downtime in that period.

The meta-lesson I took: influence without authority works best when you make it easy for people to say yes. I didn\textquotesingle t ask them to trust me --- I gave them a PoC they could verify."
\end{quote}

\textbf{Why it works:}

\begin{itemize}
\tightlist
\item
  Specific timeline, scale, stakes
\item
  Your role vs. team role is clear
\item
  The A section shows multiple concrete steps and \emph{decisions} with reasoning
\item
  You acknowledged the other person\textquotesingle s perspective (key for leadership signals)
\item
  Result is quantified AND has a comparison point
\item
  Closes with a principle --- shows self-awareness and growth
\end{itemize}

\subsection{STAR anti-patterns}\label{star-anti-patterns}

\begin{longtable}[]{@{}
  >{\raggedright\arraybackslash}p{(\columnwidth - 2\tabcolsep) * \real{0.4524}}
  >{\raggedright\arraybackslash}p{(\columnwidth - 2\tabcolsep) * \real{0.5476}}@{}}
\toprule\noalign{}
\rowcolor{acc!16}
\begin{minipage}[b]{\linewidth}\raggedright
Anti-pattern
\end{minipage} & \begin{minipage}[b]{\linewidth}\raggedright
Fix
\end{minipage} \\
\midrule\noalign{}
\endhead
\bottomrule\noalign{}
\endlastfoot
"We did X" --- no individual contribution & Start sentences with "I specifically..." \\
\rowcolor{zebra}
Result: "It was a great success" & Say: "We reduced latency by X\%, saving \$Y/month" \\
Answer runs 8+ minutes & Rehearse to 3-4 minutes; pause and ask if they want more \\
\rowcolor{zebra}
Vague situation & Name the product, the quarter, the user count \\
Skipping alternatives in the Action & Say: "I considered X, Y, Z. I chose Y because..." \\
\end{longtable}

\section{Core Themes}\label{core-themes}

\subsection{Leadership}\label{leadership}

\textbf{What interviewers look for:} Did you step up without being asked? Did you set direction, align others, make hard calls? Do you distinguish management from leadership?

\textbf{Sample question:} "Tell me about a time you led a project or initiative."

\textbf{Strong answer demonstrates:}

\begin{itemize}
\tightlist
\item
  You identified the need (not just responded to an assignment)
\item
  You defined scope and priorities --- and \emph{cut} things
\item
  You drove alignment across people who didn\textquotesingle t report to you
\item
  You stayed accountable when things went wrong
\end{itemize}

\textbf{Interview takeaway: Leadership = initiating + aligning + delivering + owning outcomes.}

\subsection{Ownership}\label{ownership}

\textbf{What interviewers look for:} Do you treat problems as yours to solve, or do you pass them up the chain? Do you follow through? Do you own failures?

\textbf{Sample question:} "Describe a time you saw a problem outside your immediate scope and took ownership of it."

\textbf{Strong answer demonstrates:}

\begin{itemize}
\tightlist
\item
  You didn\textquotesingle t wait for permission
\item
  You kept stakeholders informed rather than going dark
\item
  You saw it through --- you didn\textquotesingle t "raise the flag and hand it off"
\item
  You would do the same thing again (not resentful)
\end{itemize}

\textbf{Interview takeaway: Ownership = "my problem until it\textquotesingle s solved, regardless of org chart."}

\subsection{Conflict Resolution}\label{conflict-resolution}

\textbf{What interviewers look for:} Do you avoid conflict (bad) or escalate too quickly (bad)? Can you disagree productively and maintain relationships?

\textbf{Sample question:} "Tell me about a time you disagreed with a colleague or manager. How did you handle it?"

\textbf{Strong answer demonstrates:}

\begin{itemize}
\tightlist
\item
  You engaged directly, not passive-aggressively
\item
  You listened first, then advocated
\item
  You used data or structured argument, not emotion
\item
  You reached a resolution (or committed even if unresolved --- see Disagree-and-Commit)
\item
  The relationship survived and ideally improved
\end{itemize}

\textbf{Interview takeaway: Good conflict = direct + data-driven + resolution-focused. Bad conflict = avoided or escalated.}

\subsection{Cross-team Collaboration}\label{cross-team-collaboration}

\textbf{What interviewers look for:} Can you work across org boundaries, different incentives, different timelines? Do you create mutual wins or "win" at others\textquotesingle{} expense?

\textbf{Sample question:} "Give me an example of a time you worked with a team that had different goals than yours."

\textbf{Strong answer demonstrates:}

\begin{itemize}
\tightlist
\item
  You proactively mapped their incentives before asking for help
\item
  You found shared goals or traded value
\item
  You maintained the relationship long-term
\item
  Result benefited both teams, not just yours
\end{itemize}

\textbf{Interview takeaway: Cross-team = understand their incentives, create mutual wins, maintain relationships.}

\subsection{Ambiguity}\label{ambiguity}

\textbf{What interviewers look for:} When requirements are unclear, do you freeze? Do you ask good questions? Do you make a reasonable bet, execute, and correct course?

\textbf{Sample question:} "Tell me about a time you had to make a decision with incomplete information."

\textbf{Strong answer demonstrates:}

\begin{itemize}
\tightlist
\item
  You identified what you knew vs. didn\textquotesingle t know
\item
  You made a time-boxed bet with explicit assumptions
\item
  You built in checkpoints to validate assumptions early
\item
  You communicated uncertainty to stakeholders proactively
\item
  You course-corrected gracefully
\end{itemize}

\textbf{Interview takeaway: Ambiguity = clarify \(\rightarrow\) bet \(\rightarrow\) validate \(\rightarrow\) adjust. Don\textquotesingle t wait for perfect info.}

\subsection{Mentoring}\label{mentoring}

\textbf{What interviewers look for:} Do you make the team better? Do you invest in others? Do you create leverage beyond your own output?

\textbf{Sample question:} "Tell me about a time you helped someone grow or improve."

\textbf{Strong answer demonstrates:}

\begin{itemize}
\tightlist
\item
  You identified the person\textquotesingle s specific growth area (not generic "gave feedback")
\item
  You invested time: pair programming, structured 1:1s, code reviews with explanation
\item
  You saw measurable improvement in their work
\item
  They succeeded and you feel proud, not territorial
\end{itemize}

\textbf{Interview takeaway: Mentoring = specific investment, measurable growth, you multiplied the team\textquotesingle s output.}

\subsection{Project Impact}\label{project-impact}

\textbf{What interviewers look for:} Can you connect your technical work to business outcomes? Do you measure your impact or just ship code?

\textbf{Sample question:} "Tell me about the project you\textquotesingle re most proud of and its impact."

\textbf{Strong answer demonstrates:}

\begin{itemize}
\tightlist
\item
  Business metric moved, not just "we launched it"
\item
  You understand the mechanism: your tech \(\rightarrow\) user behavior \(\rightarrow\) metric
\item
  You had a hypothesis before and validated after
\item
  You can speak to what would have happened without the project
\end{itemize}

\textbf{Interview takeaway: Impact = metric moved \(\times\) scale \(\times\) counterfactual.}

\section{Amazon\textquotesingle s 16 Leadership Principles}\label{amazons-16-leadership-principles}

\begin{longtable}[]{@{}
  >{\raggedright\arraybackslash}p{(\columnwidth - 6\tabcolsep) * \real{0.0600}}
  >{\raggedright\arraybackslash}p{(\columnwidth - 6\tabcolsep) * \real{0.2673}}
  >{\raggedright\arraybackslash}p{(\columnwidth - 6\tabcolsep) * \real{0.2759}}
  >{\raggedright\arraybackslash}p{(\columnwidth - 6\tabcolsep) * \real{0.4313}}@{}}
\toprule\noalign{}
\rowcolor{acc!16}
\begin{minipage}[b]{\linewidth}\raggedright
\#
\end{minipage} & \begin{minipage}[b]{\linewidth}\raggedright
Principle
\end{minipage} & \begin{minipage}[b]{\linewidth}\raggedright
Core Meaning
\end{minipage} & \begin{minipage}[b]{\linewidth}\raggedright
Sample Interview Question
\end{minipage} \\
\midrule\noalign{}
\endhead
\bottomrule\noalign{}
\endlastfoot
1 & \textbf{Customer Obsession} & Start with the customer and work backward & "Tell me about a time you sacrificed short-term engineering convenience for the customer experience." \\
\rowcolor{zebra}
2 & \textbf{Ownership} & Think like an owner; don\textquotesingle t pass the buck & "Describe a time you took on a problem outside your role." \\
3 & \textbf{Invent and Simplify} & Innovation + cut complexity relentlessly & "Tell me about a time you simplified a complex process." \\
\rowcolor{zebra}
4 & \textbf{Are Right, A Lot} & Strong judgment; seek diverse views & "Tell me about a time you had to make a decision with conflicting information." \\
5 & \textbf{Learn and Be Curious} & Constant self-improvement & "What\textquotesingle s the most important thing you\textquotesingle ve taught yourself in the last year?" \\
\rowcolor{zebra}
6 & \textbf{Hire and Develop the Best} & Raise the bar; invest in people & "Tell me about someone you hired or mentored and how you developed them." \\
7 & \textbf{Insist on the Highest Standards} & Never settle for "good enough" & "Tell me about a time you raised quality standards on a team." \\
\rowcolor{zebra}
8 & \textbf{Think Big} & Bold directions, not incremental & "Tell me about a time your vision was bigger than what the team initially believed." \\
9 & \textbf{Bias for Action} & Speed matters; calculated risks & "Tell me about a time you made a quick decision with imperfect information." \\
\rowcolor{zebra}
10 & \textbf{Frugality} & Do more with less; constraints breed creativity & "Tell me about a time you delivered impact under resource constraints." \\
11 & \textbf{Earn Trust} & Listen, be vulnerable, follow through & "Tell me about a time you had to rebuild trust with a colleague or stakeholder." \\
\rowcolor{zebra}
12 & \textbf{Dive Deep} & Stay connected to details; audit & "Tell me about a time you uncovered a problem by looking at the details others missed." \\
13 & \textbf{Have Backbone; Disagree and Commit} & Challenge respectfully; then commit fully & "Tell me about a time you disagreed with your manager and how it ended." \\
\rowcolor{zebra}
14 & \textbf{Deliver Results} & Outcomes matter despite obstacles & "Tell me about a time you delivered under difficult circumstances." \\
15 & \textbf{Strive to be Earth\textquotesingle s Best Employer} & Empathy + growth for your team & "Tell me about a time you supported a struggling team member." \\
\rowcolor{zebra}
16 & \textbf{Success and Scale Bring Broad Responsibility} & Ethical scale, societal impact & "Tell me about a time you made a decision that was right for users but costly to your team." \\
\end{longtable}

\textbf{Mapping tips:}

\begin{itemize}
\tightlist
\item
  LP 2 (Ownership) + LP 14 (Deliver Results) are probed in almost every Amazon loop
\item
  LP 13 (Backbone) is the one most candidates avoid preparing --- prepare 2 stories for it
\item
  LP 6 (Hire and Develop) is critical for senior+ roles; junior candidates can use mentoring stories
\end{itemize}

\section{Building Your Story Bank}\label{building-your-story-bank}

\subsection{Why a story bank}\label{why-a-story-bank}

You have \textasciitilde6-8 strong career stories. Each story, told well, can answer 4-6 different question types. Without a bank, you repeat yourself, blank on questions, or give unrelated examples. With a bank, you map questions to stories in 5 seconds.

\subsection{The matrix template}\label{the-matrix-template}

Copy this template and fill in your real stories.

\setcodelang{Python}

\begin{Shaded}
\begin{Highlighting}[]
\NormalTok{STORY BANK MATRIX}
\OperatorTok{=================}

           \OperatorTok{|}\NormalTok{ Lead }\OperatorTok{|}\NormalTok{ Own  }\OperatorTok{|}\NormalTok{ Conflict }\OperatorTok{|}\NormalTok{ Xteam }\OperatorTok{|}\NormalTok{ Ambig }\OperatorTok{|}\NormalTok{ Mentor }\OperatorTok{|}\NormalTok{ Impact }\OperatorTok{|}\NormalTok{ Fail  }\OperatorTok{|}\NormalTok{ LP specific}
\OperatorTok{{-}{-}{-}{-}{-}{-}{-}{-}{-}{-}{-}+{-}{-}{-}{-}{-}{-}+{-}{-}{-}{-}{-}{-}+{-}{-}{-}{-}{-}{-}{-}{-}{-}{-}+{-}{-}{-}{-}{-}{-}{-}+{-}{-}{-}{-}{-}{-}{-}+{-}{-}{-}{-}{-}{-}{-}{-}+{-}{-}{-}{-}{-}{-}{-}{-}+{-}{-}{-}{-}{-}{-}{-}+{-}{-}{-}{-}{-}{-}{-}{-}{-}{-}{-}{-}}
\NormalTok{Story A    }\OperatorTok{|}\NormalTok{  ✓   }\OperatorTok{|}\NormalTok{  ✓   }\OperatorTok{|}          \OperatorTok{|}\NormalTok{  ✓    }\OperatorTok{|}       \OperatorTok{|}        \OperatorTok{|}\NormalTok{  ✓     }\OperatorTok{|}       \OperatorTok{|}\NormalTok{ LP2, LP14}
\NormalTok{(Project X)}\OperatorTok{|}      \OperatorTok{|}      \OperatorTok{|}          \OperatorTok{|}       \OperatorTok{|}       \OperatorTok{|}        \OperatorTok{|}        \OperatorTok{|}       \OperatorTok{|}
\OperatorTok{{-}{-}{-}{-}{-}{-}{-}{-}{-}{-}{-}+{-}{-}{-}{-}{-}{-}+{-}{-}{-}{-}{-}{-}+{-}{-}{-}{-}{-}{-}{-}{-}{-}{-}+{-}{-}{-}{-}{-}{-}{-}+{-}{-}{-}{-}{-}{-}{-}+{-}{-}{-}{-}{-}{-}{-}{-}+{-}{-}{-}{-}{-}{-}{-}{-}+{-}{-}{-}{-}{-}{-}{-}+{-}{-}{-}{-}{-}{-}{-}{-}{-}{-}{-}{-}}
\NormalTok{Story B    }\OperatorTok{|}      \OperatorTok{|}      \OperatorTok{|}\NormalTok{  ✓       }\OperatorTok{|}       \OperatorTok{|}\NormalTok{  ✓    }\OperatorTok{|}        \OperatorTok{|}        \OperatorTok{|}\NormalTok{  ✓    }\OperatorTok{|}\NormalTok{ LP13, LP9}
\NormalTok{(Conflict) }\OperatorTok{|}      \OperatorTok{|}      \OperatorTok{|}          \OperatorTok{|}       \OperatorTok{|}       \OperatorTok{|}        \OperatorTok{|}        \OperatorTok{|}       \OperatorTok{|}
\OperatorTok{{-}{-}{-}{-}{-}{-}{-}{-}{-}{-}{-}+{-}{-}{-}{-}{-}{-}+{-}{-}{-}{-}{-}{-}+{-}{-}{-}{-}{-}{-}{-}{-}{-}{-}+{-}{-}{-}{-}{-}{-}{-}+{-}{-}{-}{-}{-}{-}{-}+{-}{-}{-}{-}{-}{-}{-}{-}+{-}{-}{-}{-}{-}{-}{-}{-}+{-}{-}{-}{-}{-}{-}{-}+{-}{-}{-}{-}{-}{-}{-}{-}{-}{-}{-}{-}}
\NormalTok{Story C    }\OperatorTok{|}      \OperatorTok{|}\NormalTok{  ✓   }\OperatorTok{|}          \OperatorTok{|}       \OperatorTok{|}\NormalTok{  ✓    }\OperatorTok{|}        \OperatorTok{|}\NormalTok{  ✓     }\OperatorTok{|}       \OperatorTok{|}\NormalTok{ LP3, LP9}
\NormalTok{(Ambig)    }\OperatorTok{|}      \OperatorTok{|}      \OperatorTok{|}          \OperatorTok{|}       \OperatorTok{|}       \OperatorTok{|}        \OperatorTok{|}        \OperatorTok{|}       \OperatorTok{|}
\OperatorTok{{-}{-}{-}{-}{-}{-}{-}{-}{-}{-}{-}+{-}{-}{-}{-}{-}{-}+{-}{-}{-}{-}{-}{-}+{-}{-}{-}{-}{-}{-}{-}{-}{-}{-}+{-}{-}{-}{-}{-}{-}{-}+{-}{-}{-}{-}{-}{-}{-}+{-}{-}{-}{-}{-}{-}{-}{-}+{-}{-}{-}{-}{-}{-}{-}{-}+{-}{-}{-}{-}{-}{-}{-}+{-}{-}{-}{-}{-}{-}{-}{-}{-}{-}{-}{-}}
\NormalTok{Story D    }\OperatorTok{|}      \OperatorTok{|}      \OperatorTok{|}          \OperatorTok{|}\NormalTok{  ✓    }\OperatorTok{|}       \OperatorTok{|}\NormalTok{  ✓     }\OperatorTok{|}\NormalTok{  ✓     }\OperatorTok{|}       \OperatorTok{|}\NormalTok{ LP6, LP11}
\NormalTok{(Mentored) }\OperatorTok{|}      \OperatorTok{|}      \OperatorTok{|}          \OperatorTok{|}       \OperatorTok{|}       \OperatorTok{|}        \OperatorTok{|}        \OperatorTok{|}       \OperatorTok{|}
\OperatorTok{{-}{-}{-}{-}{-}{-}{-}{-}{-}{-}{-}+{-}{-}{-}{-}{-}{-}+{-}{-}{-}{-}{-}{-}+{-}{-}{-}{-}{-}{-}{-}{-}{-}{-}+{-}{-}{-}{-}{-}{-}{-}+{-}{-}{-}{-}{-}{-}{-}+{-}{-}{-}{-}{-}{-}{-}{-}+{-}{-}{-}{-}{-}{-}{-}{-}+{-}{-}{-}{-}{-}{-}{-}+{-}{-}{-}{-}{-}{-}{-}{-}{-}{-}{-}{-}}
\NormalTok{Story E    }\OperatorTok{|}\NormalTok{  ✓   }\OperatorTok{|}      \OperatorTok{|}\NormalTok{  ✓       }\OperatorTok{|}       \OperatorTok{|}       \OperatorTok{|}        \OperatorTok{|}        \OperatorTok{|}\NormalTok{  ✓    }\OperatorTok{|}\NormalTok{ LP13, LP8}
\NormalTok{(Failure)  }\OperatorTok{|}      \OperatorTok{|}      \OperatorTok{|}          \OperatorTok{|}       \OperatorTok{|}       \OperatorTok{|}        \OperatorTok{|}        \OperatorTok{|}       \OperatorTok{|}
\OperatorTok{{-}{-}{-}{-}{-}{-}{-}{-}{-}{-}{-}+{-}{-}{-}{-}{-}{-}+{-}{-}{-}{-}{-}{-}+{-}{-}{-}{-}{-}{-}{-}{-}{-}{-}+{-}{-}{-}{-}{-}{-}{-}+{-}{-}{-}{-}{-}{-}{-}+{-}{-}{-}{-}{-}{-}{-}{-}+{-}{-}{-}{-}{-}{-}{-}{-}+{-}{-}{-}{-}{-}{-}{-}+{-}{-}{-}{-}{-}{-}{-}{-}{-}{-}{-}{-}}
\NormalTok{Story F    }\OperatorTok{|}\NormalTok{  ✓   }\OperatorTok{|}\NormalTok{  ✓   }\OperatorTok{|}          \OperatorTok{|}\NormalTok{  ✓    }\OperatorTok{|}\NormalTok{  ✓    }\OperatorTok{|}        \OperatorTok{|}\NormalTok{  ✓     }\OperatorTok{|}       \OperatorTok{|}\NormalTok{ LP1, LP14}
\NormalTok{(Big proj) }\OperatorTok{|}      \OperatorTok{|}      \OperatorTok{|}          \OperatorTok{|}       \OperatorTok{|}       \OperatorTok{|}        \OperatorTok{|}        \OperatorTok{|}       \OperatorTok{|}
\OperatorTok{{-}{-}{-}{-}{-}{-}{-}{-}{-}{-}{-}+{-}{-}{-}{-}{-}{-}+{-}{-}{-}{-}{-}{-}+{-}{-}{-}{-}{-}{-}{-}{-}{-}{-}+{-}{-}{-}{-}{-}{-}{-}+{-}{-}{-}{-}{-}{-}{-}+{-}{-}{-}{-}{-}{-}{-}{-}+{-}{-}{-}{-}{-}{-}{-}{-}+{-}{-}{-}{-}{-}{-}{-}+{-}{-}{-}{-}{-}{-}{-}{-}{-}{-}{-}{-}}
\NormalTok{Story G    }\OperatorTok{|}      \OperatorTok{|}      \OperatorTok{|}          \OperatorTok{|}\NormalTok{  ✓    }\OperatorTok{|}       \OperatorTok{|}\NormalTok{  ✓     }\OperatorTok{|}        \OperatorTok{|}       \OperatorTok{|}\NormalTok{ LP6, LP11}
\NormalTok{(Collab)   }\OperatorTok{|}      \OperatorTok{|}      \OperatorTok{|}          \OperatorTok{|}       \OperatorTok{|}       \OperatorTok{|}        \OperatorTok{|}        \OperatorTok{|}       \OperatorTok{|}
\OperatorTok{{-}{-}{-}{-}{-}{-}{-}{-}{-}{-}{-}+{-}{-}{-}{-}{-}{-}+{-}{-}{-}{-}{-}{-}+{-}{-}{-}{-}{-}{-}{-}{-}{-}{-}+{-}{-}{-}{-}{-}{-}{-}+{-}{-}{-}{-}{-}{-}{-}+{-}{-}{-}{-}{-}{-}{-}{-}+{-}{-}{-}{-}{-}{-}{-}{-}+{-}{-}{-}{-}{-}{-}{-}+{-}{-}{-}{-}{-}{-}{-}{-}{-}{-}{-}{-}}
\NormalTok{Story H    }\OperatorTok{|}      \OperatorTok{|}\NormalTok{  ✓   }\OperatorTok{|}          \OperatorTok{|}       \OperatorTok{|}       \OperatorTok{|}        \OperatorTok{|}        \OperatorTok{|}\NormalTok{  ✓    }\OperatorTok{|}\NormalTok{ LP2, LP5}
\NormalTok{(Screwup)  }\OperatorTok{|}      \OperatorTok{|}      \OperatorTok{|}          \OperatorTok{|}       \OperatorTok{|}       \OperatorTok{|}        \OperatorTok{|}        \OperatorTok{|}       \OperatorTok{|}
\end{Highlighting}
\end{Shaded}

\subsection{Story bank rules}\label{story-bank-rules}

\begin{enumerate}
\def\labelenumi{\arabic{enumi}.}
\tightlist
\item
  \textbf{6-8 stories is the target.} More becomes unmanageable. Fewer leaves gaps.
\item
  \textbf{Each story must cover at least 2 themes} to be worth memorizing.
\item
  \textbf{Include at least one genuine failure story} --- interviewers can always tell when a "failure" is really a success in disguise.
\item
  \textbf{Include at least one conflict story} where you disagreed with a manager or senior stakeholder.
\item
  \textbf{Include one mentoring story} for every senior+ role.
\item
  \textbf{Have your most impactful project story} ready with 3 numbers: scale, improvement, business outcome.
\item
  \textbf{Vary recency.} Don\textquotesingle t use all stories from one job. Spread across at least 2 roles.
\end{enumerate}

\subsection{Story template (fill in for each story)}\label{story-template-fill-in-for-each-story}

\setcodelang{}

\begin{verbatim}
Story Name: ____________________
Company/Role: ___________________
Year: __________________________

S — Situation (2 sentences max):


T — My specific task/role:


A — Key actions I took (bullet list of 4-6):
• 
• 
• 
• 

R — Result (must have at least 2 numbers):


What it demonstrates:
☐ Leadership      ☐ Ownership       ☐ Conflict Resolution
☐ Cross-team      ☐ Ambiguity       ☐ Mentoring
☐ Impact          ☐ Failure/Learn   ☐ Bias for Action

LP coverage: ____________________
\end{verbatim}

\section{Architecture / Frameworks}\label{architecture--frameworks}

\subsection{STAR structure (ASCII)}\label{star-structure-ascii}

\setcodelang{}

\begin{verbatim}
┌──────────────────────────────────────────────────────┐
│                  STAR ANSWER STRUCTURE               │
├─────────────┬────────────────────────────────────────┤
│ SITUATION   │ Brief. What was the context?           │
│  (10%)      │ When, where, what was at stake?        │
├─────────────┼────────────────────────────────────────┤
│ TASK        │ What was YOUR specific job here?       │
│  (10%)      │ (Not the team's. Yours.)               │
├─────────────┼────────────────────────────────────────┤
│ ACTION      │ What did YOU do? Step by step.         │
│  (60%)      │ WHY did you choose this? Alternatives? │
│             │ How did you influence/decide/unblock?  │
├─────────────┼────────────────────────────────────────┤
│ RESULT      │ What happened? NUMBERS.                │
│  (20%)      │ What would NOT have happened w/o you? │
│             │ What did you learn?                    │
└─────────────┴────────────────────────────────────────┘
\end{verbatim}

\subsection{Story matrix grid (themes vs. LPs)}\label{story-matrix-grid-themes-vs-lps}

\setcodelang{}

\begin{verbatim}
STORY → THEME MAPPING GRID

Stories:   A    B    C    D    E    F    G    H
           │    │    │    │    │    │    │    │
Lead   ────┼────┼────┼────┼────┼────┼────┼────┤
Own    ────┼────┼────┼────┼────┼────┼────┼────┤
Confli ────┼────┼────┼────┼────┼────┼────┼────┤
Xteam  ────┼────┼────┼────┼────┼────┼────┼────┤
Ambig  ────┼────┼────┼────┼────┼────┼────┼────┤
Mentor ────┼────┼────┼────┼────┼────┼────┼────┤
Impact ────┼────┼────┼────┼────┼────┼────┼────┤
Fail   ────┼────┼────┼────┼────┼────┼────┼────┤

Goal: No column has fewer than 2 checkmarks.
      No row has only 1 checkmark.
\end{verbatim}

\subsection{Scoring rubric (how interviewers actually score)}\label{scoring-rubric-how-interviewers-actually-score}

\setcodelang{}

\begin{verbatim}
BEHAVIORAL SCORING RUBRIC (4-point scale)

4 — STRONG HIRE
    • Specific, detailed, clearly their own action
    • Multiple data points / metrics in result
    • Demonstrates exactly the LP/theme being probed
    • Shows learning or growth
    • Appropriate scope for the level being hired

3 — HIRE
    • Specific situation, clear action
    • Some metrics but could be stronger
    • Demonstrates theme adequately
    • Appropriate level scope

2 — NO HIRE
    • Vague situation or shared action ("we did")
    • No metrics; "it went well"
    • Indirectly related to theme
    • Scope seems low for level

1 — STRONG NO HIRE
    • Can't produce a relevant example
    • Blames others for outcomes
    • No learning from failures
    • Contradicts the LP (e.g., "customer feedback doesn't matter to me")

LEVEL CALIBRATION:
Senior Engineer → expects cross-team influence, measurable impact
Staff/Principal → expects org-wide influence, sets direction, drives strategy
\end{verbatim}

\section{Real-World Example Answers}\label{real-world-example-answers}

\subsection{Example 1: Conflict resolution}\label{example-1-conflict-resolution}

\textbf{Q: "Tell me about a time you had a conflict with a teammate."}

\begin{quote}
In early Q2 two years ago, I was a senior engineer on a team that owned our search ranking pipeline. A peer engineer --- also senior --- and I had a fundamental disagreement about how to handle a major infrastructure migration. She wanted to migrate incrementally, one ranking feature at a time, over three months. I thought that was the wrong approach for our system: our features had tight coupling through a shared feature store, and migrating them independently would create a six-week window of inconsistent state that could silently degrade ranking quality.

My task was to get us aligned on an approach before we started building, because diverging mid-migration would be extremely costly.

I started by having a 1:1 with her, specifically to understand her reasoning --- not to argue. She was worried about risk: a big-bang migration would put the entire pipeline at risk in one deployment. That was a legitimate concern I hadn\textquotesingle t weighted enough. So instead of pushing back on her approach, I drafted a third option: a parallel-run migration where we built the new system alongside the old, ran both in shadow mode, and compared outputs before switching traffic. It addressed her risk concern (we could roll back at any point) and my correctness concern (both systems ran the same data simultaneously, so we\textquotesingle d catch divergence immediately).

I wrote a short doc --- two pages --- presenting all three options with an explicit risk/effort/timeline tradeoff table. I shared it with her first, incorporated her edits, and then we presented it jointly to the team lead. We aligned on the parallel-run approach.

The migration took ten weeks, slightly longer than either original estimate, but we caught two subtle divergences in shadow mode that would have caused silent ranking degradation if we\textquotesingle d gone incremental. Our ranking quality metrics (NDCG) held flat through the migration --- no regression. My peer and I ended up working together on two more projects that year; the conflict actually built more trust than if we\textquotesingle d never disagreed.
\end{quote}

\textbf{Why it works:} Specific timeline, legitimate two-sided conflict, solution that integrated both perspectives, quantified result, relationship outcome.

\subsection{Example 2: Failure}\label{example-2-failure}

\textbf{Q: "Tell me about a time you failed."}

\begin{quote}
This was about three years ago. I was leading the backend for a new real-time bidding feature that would let advertisers on our platform set dynamic floor prices. It was a high-visibility project --- the PM had committed a launch date to the VP eight weeks out.

I made a mistake in how I estimated the work. I assumed we could reuse our existing auction system with a thin adapter layer. I didn\textquotesingle t do a proper technical spike first --- I let schedule pressure push me to start building before I\textquotesingle d validated that assumption. Four weeks in, my engineer hit a fundamental incompatibility: our existing auction system used batch evaluation, and the new feature required sub-10ms per-bid decisions. The adapter approach would never work.

We were four weeks from launch with a six-week rebuild ahead of us.

I called it immediately --- same day my engineer flagged it. I wrote a clear status update that said: "Here is the problem, here is why my earlier estimate was wrong, here are three options with different timelines and scope." I didn\textquotesingle t soften it or try to make it sound recoverable if it wasn\textquotesingle t. The PM and I took Option 2 --- a reduced-scope launch that kept the new pricing logic but used a simplified synchronous path that could ship in four weeks, with a plan to add the real-time optimization in the following quarter.

We shipped on the original date with the reduced scope. The reduced version drove an 8\% increase in advertiser bids, which was actually 60\% of the projected value of the full feature. The full version shipped the following quarter.

What I changed after that: I now require a technical spike for any project that involves integrating with existing infrastructure. Even one week of upfront investigation has saved us from three similar situations in the two years since. I also tell my team: "Bad news does not age well --- surface it fast."
\end{quote}

\textbf{Why it works:} Real failure (not disguised success), takes personal accountability, result still had business impact, clear behavioral change afterward.

\subsection{Example 3: Leadership / impact}\label{example-3-leadership--impact}

\textbf{Q: "Tell me about your most impactful project."}

\begin{quote}
About two years ago, I led a project we called "cold-start rescue" on our recommendations team. The problem: new users on our platform had essentially no personalization for their first 20-30 sessions. We showed them popularity-ranked content. Our data showed that users who didn\textquotesingle t receive relevant recommendations in their first week churned at 3x the rate of users who did. With 400,000 new users per month, this was a significant retention problem.

No one had formally owned this problem. I saw it in our funnel analysis and proposed it as a project. My manager gave me one engineer and one data scientist.

My first action was to actually talk to users --- we ran 12 user interviews with people who had recently churned. The dominant signal was: "I couldn\textquotesingle t find anything I wanted to watch." We also analyzed the behavioral data and found that users who completed our onboarding survey (only 15\% did) had significantly better day-7 retention. But the survey was optional and buried.

I redesigned the strategy: instead of better algorithms on sparse data, we\textquotesingle d improve signal collection in the first session. I worked with the PM (different team) to redesign onboarding --- we made the taste-selection flow interactive and fun instead of a survey, and integrated it into the first session rather than asking upfront. I worked with the ML engineer to build a two-phase model: onboarding signals for the first 30 sessions, then our standard collaborative filtering after.

I also had to negotiate with the growth team, who owned onboarding, because my proposed changes affected their funnel metrics. I spent two weeks aligning on a shared metric --- "7-day active return rate" --- that both teams cared about, and got their buy-in to run a joint experiment.

We shipped the new onboarding to 20\% of new users. At 8 weeks: 7-day retention for new users increased 22\%. Churn in week 1 dropped 17\%. Annualized, this was estimated to retain \textasciitilde85,000 additional users per year at our average LTV of \$48/user --- roughly \$4M in annual retained revenue. The project became a model for how we run cross-team initiatives; the shared-metric framework I built has been adopted by two other team pairs.
\end{quote}

\textbf{Why it works:} Started with problem identification (not assignment), cross-team collaboration, user research, specific metrics, business impact quantified, lasting process contribution.

\section{Real-Life Analogies}\label{real-life-analogies}

\emph{One team, one season --- every trait is a play, a role, or a moment on the same field.}

\begin{longtable}[]{@{}
  >{\raggedright\arraybackslash}p{(\columnwidth - 2\tabcolsep) * \real{0.1686}}
  >{\raggedright\arraybackslash}p{(\columnwidth - 2\tabcolsep) * \real{0.8314}}@{}}
\toprule\noalign{}
\rowcolor{acc!16}
\begin{minipage}[b]{\linewidth}\raggedright
Principle
\end{minipage} & \begin{minipage}[b]{\linewidth}\raggedright
Analogy
\end{minipage} \\
\midrule\noalign{}
\endhead
\bottomrule\noalign{}
\endlastfoot
STAR Method & A highlight reel of one play: the game situation, your role on the field, the move you made, and the score it changed --- forty-five minutes of match compressed into ninety seconds of tape that shows exactly what you did. \\
\rowcolor{zebra}
Ownership & The captain who steps up to the microphone after a 3-0 loss and says "that was on me, not the goalkeeper" --- no finger-pointing, no asterisks, just full accountability for what happened on the pitch. \\
Conflict Resolution & Two midfielders arguing opposite tactics in the huddle at half-time --- each makes their case hard, the coach decides, and both players sprint back out and execute the same play with full commitment. Disagree in the locker room, commit on the field. \\
\rowcolor{zebra}
Cross-team Collaboration & Running a set piece that only works when the forwards, the wide players, the physio who cleared the striker to start, and the video analyst who spotted the defensive gap all do their job --- not just your position, the whole squad. \\
Ambiguity & The opponent switches from a back-four to a three-man defence in the 55th minute without warning. There is no time for a coaching call: you read the formation, signal your line-mates, and adapt the press before the next goal kick lands. \\
\rowcolor{zebra}
Mentoring & The club\textquotesingle s ten-year veteran who stays on the training pitch after everyone else has gone in, walking the rookie through defensive transitions one repetition at a time until the movement becomes automatic. \\
Project Impact & The season stats that actually won the title --- not the flashy individual dribbles, but the goals, assists, and clean sheets that show up in the standings and land the trophy in May. \\
\rowcolor{zebra}
Influence without Authority & The player with no armband whom every teammate looks to when the game gets tense --- not appointed, just trusted, because her reading of the game has been right too many times to ignore. \\
Leadership & Spotting that the press is breaking down before the coach sees it from the sideline, calling the defensive shape out loud mid-run, and getting eleven people moving as one without a whistle being blown. \\
\rowcolor{zebra}
Bias for Action & Taking the shot on goal in the 88th minute instead of waiting for a more perfect angle that may never come --- a 60\% chance you take now beats a 90\% chance that arrives after the final whistle. \\
Customer Obsession & Building the attack around what the fans in the stands came to see --- relentless movement, goals --- not around tactics that are theoretically sound but empty the stadium by half-time. \\
\rowcolor{zebra}
Failure / Learning & Watching the match film the morning after a defeat, pausing on every mistake you made personally, and running the corrected movement pattern in training until your body forgets the bad version --- accountability plus repetition, not excuses. \\
Quantifying Impact & Posting the actual scoreboard --- goals, assists, clean sheets, league position --- not just telling the board "we played well." The scoreboard is what gets your contract renewed. \\
\end{longtable}

\section{Memory Tricks / Mnemonics}\label{memory-tricks--mnemonics}

\subsection{Remembering STAR}\label{remembering-star}

\textbf{"Stories Tell Actual Results"}

\begin{itemize}
\tightlist
\item
  \textbf{S}tories = Situation
\item
  \textbf{T}ell = Task
\item
  \textbf{A}ctual = Action
\item
  \textbf{R}esults = Result
\end{itemize}

Or the spatial trick: \textbf{S}cene \(\rightarrow\) \textbf{T}rigger \(\rightarrow\) \textbf{A}ction \(\rightarrow\) \textbf{A}ftermath (STAA if you need a variant).

\subsection{Remembering Amazon\textquotesingle s 16 LPs (grouped)}\label{remembering-amazons-16-lps-grouped}

Group them into 4 buckets of 4:

\textbf{CUSTOMER (what you build for):}

\begin{enumerate}
\def\labelenumi{\arabic{enumi}.}
\tightlist
\item
  Customer Obsession
\item
  Ownership
\item
  Think Big
\item
  Invent and Simplify
\end{enumerate}

\textbf{QUALITY (how you build):}
5. Are Right, A Lot
6. Dive Deep
7. Insist on Highest Standards
8. Deliver Results

\textbf{PEOPLE (how you lead):}
9. Hire and Develop the Best
10. Earn Trust
11. Learn and Be Curious
12. Strive to be Earth\textquotesingle s Best Employer

\textbf{JUDGMENT (how you decide):}
13. Bias for Action
14. Have Backbone; Disagree and Commit
15. Frugality
16. Success and Scale Bring Broad Responsibility

\textbf{Memory hook for the 4 buckets:} \textbf{"CQPJ"} --- "Customers Quietly Pay Justly"

\subsection{Remembering not to say "we"}\label{remembering-not-to-say-we}

\textbf{The "I" audit:} After every practice answer, read it back and underline every "I." If there are fewer than 5 in the Action section, you\textquotesingle re probably hiding behind "we." Interviewers can\textquotesingle t score what they can\textquotesingle t see.

\subsection{The 3-number rule for results}\label{the-3-number-rule-for-results}

Every result should have: \textbf{a baseline, a change, and a scale.}

\begin{itemize}
\tightlist
\item
  Example: "Latency was 800ms (baseline), we cut it to 120ms (change), for a service handling 50M requests/day (scale)."
\item
  Mnemonic: \textbf{B-C-S} = Before \(\rightarrow\) Change \(\rightarrow\) Scale
\end{itemize}

\subsection{Pre-interview 5-minute mental warm-up}\label{pre-interview-5-minute-mental-warm-up}

\begin{enumerate}
\def\labelenumi{\arabic{enumi}.}
\tightlist
\item
  Name your 6-8 stories out loud (30 seconds)
\item
  Say the 4 STAR words + what each means (15 seconds)
\item
  State your top 3 results with their numbers (1 minute)
\item
  Take 3 slow breaths (30 seconds)
\item
  Remind yourself: "I am describing my actual experience. There\textquotesingle s no trick answer." (15 seconds)
\end{enumerate}

\section{Common Interview Questions}\label{common-interview-questions}

\subsection{Leadership}\label{leadership-1}

\begin{longtable}[]{@{}
  >{\raggedright\arraybackslash}p{(\columnwidth - 2\tabcolsep) * \real{0.5352}}
  >{\raggedright\arraybackslash}p{(\columnwidth - 2\tabcolsep) * \real{0.4648}}@{}}
\toprule\noalign{}
\rowcolor{acc!16}
\begin{minipage}[b]{\linewidth}\raggedright
Question
\end{minipage} & \begin{minipage}[b]{\linewidth}\raggedright
What it\textquotesingle s really probing
\end{minipage} \\
\midrule\noalign{}
\endhead
\bottomrule\noalign{}
\endlastfoot
"Tell me about a time you led a team or project." & Can you set direction and align without authority? \\
\rowcolor{zebra}
"How do you handle a project where the goal is unclear?" & Ambiguity tolerance + structured thinking \\
"Tell me about a time you had to prioritize ruthlessly." & Can you cut? Do you understand tradeoffs? \\
\rowcolor{zebra}
"Describe a time you changed the direction of a project mid-way." & Adaptability + decision-making \\
"When have you had to make a decision that was unpopular?" & Backbone; willing to do the hard right thing \\
\end{longtable}

\subsection{Conflict}\label{conflict}

\begin{longtable}[]{@{}
  >{\raggedright\arraybackslash}p{(\columnwidth - 2\tabcolsep) * \real{0.5572}}
  >{\raggedright\arraybackslash}p{(\columnwidth - 2\tabcolsep) * \real{0.4428}}@{}}
\toprule\noalign{}
\rowcolor{acc!16}
\begin{minipage}[b]{\linewidth}\raggedright
Question
\end{minipage} & \begin{minipage}[b]{\linewidth}\raggedright
What it\textquotesingle s really probing
\end{minipage} \\
\midrule\noalign{}
\endhead
\bottomrule\noalign{}
\endlastfoot
"Tell me about a disagreement with a peer." & Direct, constructive conflict style \\
\rowcolor{zebra}
"Tell me about a time you disagreed with your manager." & LP 13: Backbone; didn\textquotesingle t just comply silently \\
"Have you ever had to work with a difficult person?" & Emotional intelligence + resolution \\
\rowcolor{zebra}
"Tell me about a time you pushed back on a stakeholder." & Influencing up without being insubordinate \\
"Describe a time you and a colleague couldn\textquotesingle t agree. What happened?" & Do you leave conflicts unresolved? \\
\end{longtable}

\subsection{Failure / Learning}\label{failure--learning}

\begin{longtable}[]{@{}
  >{\raggedright\arraybackslash}p{(\columnwidth - 2\tabcolsep) * \real{0.5271}}
  >{\raggedright\arraybackslash}p{(\columnwidth - 2\tabcolsep) * \real{0.4729}}@{}}
\toprule\noalign{}
\rowcolor{acc!16}
\begin{minipage}[b]{\linewidth}\raggedright
Question
\end{minipage} & \begin{minipage}[b]{\linewidth}\raggedright
What it\textquotesingle s really probing
\end{minipage} \\
\midrule\noalign{}
\endhead
\bottomrule\noalign{}
\endlastfoot
"Tell me about a time you failed." & Accountability + growth mindset \\
\rowcolor{zebra}
"Tell me about your biggest mistake." & Same --- but note "biggest" --- don\textquotesingle t give a trivial example \\
"What would you do differently if you could repeat a project?" & Self-awareness \\
\rowcolor{zebra}
"Tell me about a time your project didn\textquotesingle t deliver expected results." & Honest reflection + what you learned \\
"Have you ever made a decision you later regretted?" & Judgment quality + learning \\
\end{longtable}

\subsection{Ambiguity}\label{ambiguity-1}

\begin{longtable}[]{@{}
  >{\raggedright\arraybackslash}p{(\columnwidth - 2\tabcolsep) * \real{0.6282}}
  >{\raggedright\arraybackslash}p{(\columnwidth - 2\tabcolsep) * \real{0.3718}}@{}}
\toprule\noalign{}
\rowcolor{acc!16}
\begin{minipage}[b]{\linewidth}\raggedright
Question
\end{minipage} & \begin{minipage}[b]{\linewidth}\raggedright
What it\textquotesingle s really probing
\end{minipage} \\
\midrule\noalign{}
\endhead
\bottomrule\noalign{}
\endlastfoot
"Tell me about a time requirements kept changing." & Adaptability under ambiguity \\
\rowcolor{zebra}
"How do you start a project when you have little information?" & Structured approach to ambiguity \\
"Tell me about a time you made a decision with incomplete data." & Bias for action + calibrated risk \\
\rowcolor{zebra}
"Describe a time you had to figure out what to build without a clear spec." & Product thinking + initiative \\
\end{longtable}

\subsection{Impact}\label{impact}

\begin{longtable}[]{@{}
  >{\raggedright\arraybackslash}p{(\columnwidth - 2\tabcolsep) * \real{0.5562}}
  >{\raggedright\arraybackslash}p{(\columnwidth - 2\tabcolsep) * \real{0.4438}}@{}}
\toprule\noalign{}
\rowcolor{acc!16}
\begin{minipage}[b]{\linewidth}\raggedright
Question
\end{minipage} & \begin{minipage}[b]{\linewidth}\raggedright
What it\textquotesingle s really probing
\end{minipage} \\
\midrule\noalign{}
\endhead
\bottomrule\noalign{}
\endlastfoot
"What\textquotesingle s the project you\textquotesingle re most proud of?" & Impact measurement + attribution \\
\rowcolor{zebra}
"What\textquotesingle s the most technically complex thing you\textquotesingle ve built?" & Depth + ability to explain \\
"Tell me about a time your work had significant business impact." & Business awareness + impact quantification \\
\rowcolor{zebra}
"What\textquotesingle s the hardest technical problem you\textquotesingle ve solved?" & Problem-solving + communication \\
\end{longtable}

\subsection{Teamwork / Cross-team}\label{teamwork--cross-team}

\begin{longtable}[]{@{}
  >{\raggedright\arraybackslash}p{(\columnwidth - 2\tabcolsep) * \real{0.6263}}
  >{\raggedright\arraybackslash}p{(\columnwidth - 2\tabcolsep) * \real{0.3737}}@{}}
\toprule\noalign{}
\rowcolor{acc!16}
\begin{minipage}[b]{\linewidth}\raggedright
Question
\end{minipage} & \begin{minipage}[b]{\linewidth}\raggedright
What it\textquotesingle s really probing
\end{minipage} \\
\midrule\noalign{}
\endhead
\bottomrule\noalign{}
\endlastfoot
"Give an example of when you had to collaborate with another team." & Cross-functional skills \\
\rowcolor{zebra}
"Tell me about a time you had to earn buy-in from skeptical stakeholders." & Influence + communication \\
"Describe a situation where you helped someone else succeed." & Mentoring + team-first mindset \\
\rowcolor{zebra}
"How have you handled a situation where teams had conflicting priorities?" & Negotiation + shared goal-finding \\
\end{longtable}

\section{Senior/Staff-Level Discussion Points}\label{seniorstaff-level-discussion-points}

\subsection{Scope expectations by level}\label{scope-expectations-by-level}

\begin{longtable}[]{@{}
  >{\raggedright\arraybackslash}p{(\columnwidth - 2\tabcolsep) * \real{0.2350}}
  >{\raggedright\arraybackslash}p{(\columnwidth - 2\tabcolsep) * \real{0.7650}}@{}}
\toprule\noalign{}
\rowcolor{acc!16}
\begin{minipage}[b]{\linewidth}\raggedright
Level
\end{minipage} & \begin{minipage}[b]{\linewidth}\raggedright
Expected behavioral scope
\end{minipage} \\
\midrule\noalign{}
\endhead
\bottomrule\noalign{}
\endlastfoot
Mid-level (L4/SWE II) & Team-level impact; unblocking self; solid on STAR \\
\rowcolor{zebra}
Senior (L5/SWE III) & Cross-team influence; measurable business impact; leading without authority; mentoring 1-2 people \\
Staff/L6+ & Org-level impact; setting technical strategy; influencing across orgs; creating leverage through others \\
\end{longtable}

\textbf{At staff level, interviewers explicitly listen for:} org-wide scope, not just team-wide. If your stories are all "I fixed it on my team," that scores senior, not staff.

\subsection{Influence without authority}\label{influence-without-authority}

This is the defining skill of senior+ engineers. Framework:

\setcodelang{}

\begin{verbatim}
INFLUENCE WITHOUT AUTHORITY PLAYBOOK
=====================================
1. UNDERSTAND their incentives, not yours
   → "What does success look like for your team this quarter?"

2. FIND the shared goal
   → "What if we framed it as X — does that serve both teams?"

3. BUILD credibility through data first
   → Write the doc. Run the PoC. Show, don't tell.

4. MAKE it easy to say yes
   → Offer to do the work. Give them the win too.

5. ACCEPT the decision if overruled
   → Disagree-and-commit. State your view once, clearly, then execute.
\end{verbatim}

\subsection{Driving org-level impact}\label{driving-org-level-impact}

For staff+ roles, you should have stories about:

\begin{itemize}
\tightlist
\item
  Establishing a technical standard/convention others adopted
\item
  Identifying a systemic problem (not just your team\textquotesingle s) and driving the fix
\item
  Influencing a roadmap decision through technical argument, not just voting
\item
  Creating leverage: a tool, a platform, a process that accelerated other teams
\end{itemize}

\subsection{Handling disagreement with leadership}\label{handling-disagreement-with-leadership}

This is the hardest conversation and the most important to have a real story for.

\textbf{The disagree-and-commit framework:}

\setcodelang{}

\begin{verbatim}
Step 1: State your position once, clearly, with data
        ("Here is my concern and here is the evidence.")

Step 2: Make sure they heard it
        ("I want to make sure I'm being clear — my concern is X.")

Step 3: Ask to understand their decision
        ("Help me understand what I might be missing.")

Step 4: If they decide differently, commit fully
        ("I understand. I'll execute this and make it work.")

Step 5: Document your disagreement privately (your own notes)
        → If things go as you predicted, revisit professionally.

What NOT to do:
- Raise it publicly in a group setting (embarrassing)
- Go silent and passive-aggressively not-do-the-work
- Relitigate the decision at every subsequent meeting
- Tell your team "I don't agree with this but we have to do it"
\end{verbatim}

\subsection{For "influence up" stories at senior+ level}\label{for-influence-up-stories-at-senior-level}

Interviewers want to see you influenced a \emph{decision}, not just an implementation. "I convinced my manager to give us two more weeks" is not the same as "I changed the technical direction of a \$2M initiative."

\section{Typical Mistakes Candidates Make}\label{typical-mistakes-candidates-make}

\begin{longtable}[]{@{}
  >{\raggedright\arraybackslash}p{(\columnwidth - 4\tabcolsep) * \real{0.2778}}
  >{\raggedright\arraybackslash}p{(\columnwidth - 4\tabcolsep) * \real{0.3544}}
  >{\raggedright\arraybackslash}p{(\columnwidth - 4\tabcolsep) * \real{0.3679}}@{}}
\toprule\noalign{}
\rowcolor{acc!16}
\begin{minipage}[b]{\linewidth}\raggedright
Mistake
\end{minipage} & \begin{minipage}[b]{\linewidth}\raggedright
What happens
\end{minipage} & \begin{minipage}[b]{\linewidth}\raggedright
Fix
\end{minipage} \\
\midrule\noalign{}
\endhead
\bottomrule\noalign{}
\endlastfoot
\textbf{Saying "we" throughout} & Interviewer can\textquotesingle t score you; can\textquotesingle t tell what you did & Audit every sentence. Start with "I specifically..." \\
\rowcolor{zebra}
\textbf{No metrics in result} & Answer sounds like a story, not an outcome & Always have baseline + change + scale \\
\textbf{Choosing safe/easy failure story} & Signals lack of self-awareness or dishonesty & Pick a real failure with real stakes. Show what you learned. \\
\rowcolor{zebra}
\textbf{Rambling past 5 minutes} & Interviewer loses the thread; no time for follow-ups & Practice with a timer. 3-4 minutes target. \\
\textbf{Blaming the team or org for failures} & Signals poor ownership; red flag on LP2 & Own your part. "I could have..." not "the team didn\textquotesingle t..." \\
\rowcolor{zebra}
\textbf{Generic leadership story} & "I volunteered to run the meeting" --- too weak for senior & Match story scope to the level. Senior = cross-team+ \\
\textbf{Skipping the A section} & SIT \(\rightarrow\) RESULT with no action detail & The A section is 60\% of your score. Don\textquotesingle t rush it. \\
\rowcolor{zebra}
\textbf{No conflict or failure story prepared} & When asked, you blank or give a non-answer & Have 2 each of conflict and failure stories ready \\
\textbf{Theoretical answers} & "In that situation I would..." & Always past-tense. "Here\textquotesingle s what I did." \\
\rowcolor{zebra}
\textbf{Not adapting to the company} & Amazon LP story that never mentions customer & Map your story emphasis to the company\textquotesingle s framework \\
\end{longtable}

\section{FAANG Interview Tips}\label{faang-interview-tips}

\subsection{Amazon}\label{amazon}

\begin{itemize}
\tightlist
\item
  \textbf{Prepare 2 stories per LP} for the LPs most likely to be covered (LP2, LP13, LP14 are probed most often)
\item
  Amazon interviewers will ask follow-up questions to probe depth --- don\textquotesingle t give a shallow A section
\item
  The Bar Raiser will ask harder follow-ups than other interviewers. Expect: "What would you do differently?" and "What were the key risks you saw?"
\item
  "Think Big" (LP8) is a trap for engineers: make sure you have a story that\textquotesingle s ambitious, not just technically clever
\item
  Explicitly say "I" and "my decision" --- Amazon interviewers are trained to flag "we" answers
\item
  It\textquotesingle s okay to ask for 30 seconds to think of an example
\end{itemize}

\subsection{Google}\label{google}

\begin{itemize}
\tightlist
\item
  Google\textquotesingle s Googleyness round explicitly scores "comfort with ambiguity" and "intellectual humility" --- have stories showing you were wrong and updated
\item
  Leadership at Google = emergent. They do NOT want stories where your manager assigned you to lead. They want "I saw a gap and stepped up."
\item
  Google looks for curiosity --- often asks "what did you learn?" as a follow-up
\item
  Googleyness includes "does the right thing" --- have a story about an ethical tradeoff or doing something costly because it was right
\item
  In Google loops, your behavioral interview is reviewed by the hiring committee who also see your coding and design scores --- they look for consistency
\end{itemize}

\subsection{Meta}\label{meta}

\begin{itemize}
\tightlist
\item
  Meta values directness: if you were blunt with a colleague, that\textquotesingle s OK to say --- frame it as "I gave direct feedback" not "I told them they were wrong"
\item
  "Move fast" --- show bias for action. Stories about shipping a v1, getting feedback, and iterating fast are valued
\item
  Cross-functional stories are crucial --- Meta\textquotesingle s orgs are highly cross-functional; show you can work with PM, design, data
\item
  Meta explicitly looks for people who can operate without tight process --- show you\textquotesingle ve worked in ambiguous environments
\end{itemize}

\subsection{Microsoft}\label{microsoft}

\begin{itemize}
\tightlist
\item
  Microsoft\textquotesingle s post-Nadella culture emphasizes growth mindset --- the learning from failure matters as much as the failure itself
\item
  Empathy and inclusion language lands well: "I made sure the quieter engineers on the team had space to contribute"
\item
  Avoid stories about "I did it faster than the team" --- collaborative framing is valued more than individual heroics
\item
  "Clarity, Energy, Success" is Nadella\textquotesingle s framework --- your stories can loosely map to these
\end{itemize}

\subsection{General FAANG tips}\label{general-faang-tips}

\begin{itemize}
\tightlist
\item
  \textbf{Ask a clarifying question first} if you need to ("Can I take a moment?" or "Should I focus on a technical leadership example or a people leadership example?")
\item
  \textbf{Pause before answering} --- 10 seconds of thinking produces a better story than 3 minutes of rambling
\item
  \textbf{End strong} --- the interviewer remembers your last sentence. Make it a principle or a quantified outcome, not a trailing "...and yeah, that\textquotesingle s basically it"
\item
  \textbf{Mirror the company\textquotesingle s language} --- Amazon: "I dove deep," "I thought big." Google: "I was curious about," "I built consensus." Meta: "We shipped fast and iterated."
\item
  \textbf{Ask at the end:} "Is there anything about my answers you\textquotesingle d like me to expand on?" --- shows confidence and gives you a second chance
\end{itemize}

\section{Revision Cheat Sheet}\label{revision-cheat-sheet}

\subsection{10-minute summary}\label{10-minute-summary}

\textbf{Behavioral rounds screen for:} judgment, ownership, influence, learning. Not just technical skill.

\textbf{STAR = Situation (10\%) \(\rightarrow\) Task (10\%) \(\rightarrow\) Action (60\%) \(\rightarrow\) Result (20\%)}

\textbf{The most common failure:} Spending too much time on S and T, too little on A and R.

\textbf{Story bank:} 6-8 stories, each covering multiple themes. One failure, one conflict, one big impact, one cross-team, one mentoring (if senior+).

\textbf{Amazon:} 16 LPs. Key ones: Ownership, Deliver Results, Backbone. Always say "I." Always have data.

\textbf{Google:} Googleyness = humble + curious + does-right-thing. Leadership = emergent, not assigned.

\textbf{Numbers in every result:} baseline + change + scale. "We reduced latency 40\% for 50M daily users."

\textbf{Disagree-and-commit:} State your view once with data. Then execute fully if overruled. Don\textquotesingle t relitigate.

\textbf{Influence without authority:} Understand their incentives. Build credibility with data. Make it easy to say yes.

\subsection{Key points (top 10)}\label{key-points-top-10}

\begin{enumerate}
\def\labelenumi{\arabic{enumi}.}
\tightlist
\item
  Behavioral rounds are co-equal with coding at FAANG --- prepare equally hard
\item
  STAR Action section is 60\% of your score --- don\textquotesingle t skim it
\item
  "We" answers score 0 --- say "I" and own your specific contribution
\item
  Quantify every result: before/after + scale
\item
  Have a real failure story --- a non-failure "failure" is a red flag
\item
  Have a conflict-with-senior-stakeholder story ready
\item
  Build a story bank of 6-8 stories before the interview
\item
  Match your story scope to the level (senior = cross-team influence)
\item
  At Amazon, know all 16 LPs; practice 2 stories per most-probed LP
\item
  Disagree-and-commit: express once with data, then execute fully
\end{enumerate}

\subsection{Pre-interview checklist}\label{pre-interview-checklist}

\textbf{One week before:}

\begin{itemize}
\tightlist
\item[$\square$]
  Story bank matrix filled in (6-8 stories, covering all themes)
\item[$\square$]
  Each story written out in STAR format with 2+ numbers in result
\item[$\square$]
  Amazon LPs reviewed; top 6 LPs mapped to stories
\item[$\square$]
  Conflict story ready (peer AND manager/senior level)
\item[$\square$]
  Failure story ready (real failure, real learning)
\item[$\square$]
  Mentoring story ready (if interviewing for senior+)
\end{itemize}

\textbf{Day before:}

\begin{itemize}
\tightlist
\item[$\square$]
  Rehearse top 4 stories out loud (not in your head --- out loud)
\item[$\square$]
  Review the company\textquotesingle s values/principles (LP list, Googleyness, etc.)
\item[$\square$]
  Know your 3 numbers for your biggest impact story
\item[$\square$]
  Get 7-8 hours of sleep
\end{itemize}

\textbf{Day of:}

\begin{itemize}
\item[$\square$]
  5-minute mental warm-up (name your stories + numbers)
\item[$\square$]
  Have water nearby (dry mouth from nerves)
\item[$\square$]
  It\textquotesingle s okay to pause 10-15 seconds before answering
\item[$\square$]
  End each answer with your result number or a principle statement
\end{itemize}

\subsection{Most important concepts at a glance}\label{most-important-concepts-at-a-glance}

\setcodelang{}

\begin{verbatim}
STAR: Situation → Task → Action(60%) → Result(numbers)

Story bank: 6-8 stories × 8 themes matrix

3-number rule: Before → Change → Scale

Influence without authority:
  Understand incentives → shared goal → credibility → make it easy → commit

Disagree-and-commit:
  State once + data → understand their view → commit fully if overruled

Amazon LPs in 4 buckets:
  CUSTOMER: Obsession, Ownership, Think Big, Invent+Simplify
  QUALITY:  Are Right, Dive Deep, High Standards, Deliver Results
  PEOPLE:   Hire+Develop, Earn Trust, Learn+Curious, Best Employer
  JUDGMENT: Bias for Action, Backbone, Frugality, Broad Responsibility

Senior+ scope: cross-team influence + measurable business impact
Staff+ scope:  org-level influence + leverage through others

Never say "we" in Action section.
Never give a result without numbers.
Never give a "failure" that isn't actually a failure.
\end{verbatim}

\end{document}
